% DAFC_HowToUse
% 0709追記分
\documentclass[uplatex,dvipdfmx,a4paper,12pt]{jsreport}
%\usepackage[utf8]{inputenc}
\usepackage{makeidx}

\usepackage{amsmath}
\usepackage{amsfonts}




\begin{document}
\subsection{変数の代入}
    変数への代入は\texttt{->}または\texttt{=}で行う。
    大文字と小文字は区別して扱う。
\begin{verbatim}
    > X= 123
    > X
    [1] 123
\end{verbatim}

    Rでは値がベクトルとして扱われており、\texttt{[\ ]}内の数は要素番号を意味する。
    一般のベクトルを入力するときは\fbox{\texttt{c()}}という関数を用いる\footnote{cはcombineまたはcomcatenate(結合)の意味}

\begin{verbatim}
    > X= c(3.12, 1.41, 1.72) # 3.14, 1.41, 1.72を値として代入
    > X
    [1] 3.12, 1.41, 1.72
\end{verbatim}

    作業を終了するには\fbox{\texttt{q()}}を打ち込む。

\section{データの入力}
    少数のデータであれば\texttt{c()}を使って入力してもよい。
\begin{verbatim}
    > hight= c(168.5, 172.8, 159.0)
    > weight= c(69.5, 75.0, 56.5)
    # BMIを計算してみる、
    BMI= weight/ (hight/ 100)^ 2 # ^は指数
    [1] 24.47851, 25.11735, 22.34880
\end{verbatim}
    小数点以下を丸めるには\fbox{\texttt{round()}}関数を使い、第一引数には処理対象、第二引数には小数第何位まで表示するかを書く。
\begin{verbatim}
    > round(BMI, 1)
    [1] 24.5, 25.1, 22.3
    > round(BMI, 3)
    [1] 24.479, 25.117, 22.349
\end{verbatim}
%\begin{souremark}
    値の列はコロン\texttt{:}を使って\texttt{c(3:6)= c(3, 4, 5, 6)}のように書くこともできる。\\
    増分が一定ならば、始値、終値、増分を使って\texttt{seq(3, 9, 2)= seq(from= 3, to= 9, by= 2)= c(3, 5, 7, 9)}と書くことができる。
%\end{souremark}
%\begin{souremark}
    同じ値が繰り返されるなら\fbox{\texttt{rep()}}関数を使う。第一引数には対象、第二引数には繰り返し回数を書く。
\begin{verbatim}
    > rep(2, 5)
    [1] 2, 2, 2, 2, 2
\end{verbatim}
%\end{souremark}

\section{データフレーム}
    複数のデータを扱うときは\underline{データフレーム関数} \fbox{\texttt{deta.frame()}}を使うとよい。
\begin{verbatim}
    > drop= c(10.3, 9.7, 10.2)
    > ph= c(7.4, 7.2, 7.4)
    > X= deta.frame(drop, ph)
    > X
    [1]drop   ph
       10.3   7.4
       9.7    7.2
       10.2   7.4
\end{verbatim}
\end{document}